\chapter{Conclusiones}

Se alcanzó exitosamente el objetivo general: desarrollar un modelo de automatización basado en algoritmos genéticos para la asignación de tareas en un ERP de talleres industriales, comparándolo con el método SPT bajo condiciones reproducibles.

\textbf{Implementación del Algoritmo Genético.} Se desarrolló un AG con representación dual (asignación de operarios + prioridades de secuenciación), función de aptitud ponderada y operadores genéticos calibrados (300 generaciones, 15\% mutación, selección por ranking). El modelo es escalable, reproducible y ejecutable en tiempos prácticos.

\textbf{Desempeño Comparativo.} Aunque AG y SPT ejecutan idéntico número de tareas (29.53 promedio en 30 instancias), el AG demuestra superior estructuración: makespan 15.6\% menor (157.37 vs 186.53 minutos) y balance de carga 14.3\% mejor (desviación estándar 30.10 vs 35.11). Esta eficiencia genera valor industrial mediante cronogramas más compactos y distribución equitativa de trabajo.

\textbf{Impacto Industrial.} La reducción del 15.6\% en makespan representa 29.16 minutos diarios de eficiencia adicional, equivalente a 117.6 horas anuales (aproximadamente 15 días de productividad sin incremento de recursos). El mejor balance reduce cuellos de botella, fatiga del personal y errores operativos. Para Rectificadora Recti-Ram, esto implica mayor flexibilidad operativa, cumplimiento de plazos más confiable y capacidad para aceptar mayor volumen de órdenes.

\textbf{Validez Técnica.} El modelo formalizado proporciona marco riguroso. La validación experimental confirma: (1) Convergencia rápida en 300 generaciones; (2) Factibilidad de todas las soluciones bajo restricciones reales; (3) Estabilidad superior a SPT; (4) Reproducibilidad completa mediante semilla aleatoria fija y documentación detallada.

\textbf{Contribuciones.} El proyecto formaliza empíricamente un problema de asignación en contexto ERP real, demostrando que técnicas evolutivas superan heurísticas clásicas en este dominio. Proporciona metodología reproducible aplicable a PYMES mediante software de código abierto (Python, PyGAD), reduciendo barreras de adopción.

\textbf{Oportunidades Futuras.} Las limitaciones abren oportunidades de extensión: (1) Modelos híbridos AG-PSO o AG-ACO; (2) Job Shop Scheduling flexible con selección dinámica de máquinas; (3) Optimización multiobjetivo integrando inventarios y costos; (4) Aprendizaje automático para ajuste dinámico de parámetros; (5) Integración con IoT e Industria 4.0; (6) Computación cuántica (quantum annealing) para instancias mayores.

\textbf{Conclusión Final.} Los algoritmos genéticos demuestran ser herramienta viable para automatizar asignación de tareas industriales. Su superior desempeño respecto a SPT, adaptabilidad y reproducibilidad los posicionan como componente valiosa en sistemas de soporte de decisión operacional. Sin embargo, la implementación exitosa requiere comprensión profunda del dominio, validación rigurosa en ambiente real, diálogo continuo con usuarios finales y compromiso con mejora iterativa. La evolución natural del trabajo conduce hacia técnicas híbridas, modelos flexibles de Job Shop, integración con ERP e IoT, y arquitecturas emergentes como aprendizaje automático y computación cuántica, posicionando este aporte en la frontera de tecnologías emergentes en optimización industrial.
